 % !TeX root = book.tex

\chapter{The Prolog Web}\label{sec:the-prolog-web}

\begin{quotation}
\noindent Imagine the Web wrapped in Prolog, running on top of a distributed architecture -- a network of nodes exposing HTTP and WebSocket APIs, speaking web formats such as JSON. Think of it as a high-level Web, with Prolog agents capable of serving answers to queries -- answers that follow from what the Web knows. Moreover, imagine it being programmable, allowing Web Prolog source code to flow in either direction: from client to node, or from node to client. This is what \textbf{the Prolog Web} is all about.
\vspace{-0.2cm}
\flushright\textit{The Prolog Web -- the elevator pitch}
\end{quotation}
\vspace{0.5cm}

\noindent The previous chapters introduced Web Prolog as a language
and Prolog agents as programmable entities. This chapter turns to
the third element of the Prolog Trinity: the \emph{Prolog Web}
itself -- a network of nodes hosting actors that communicate across
node boundaries. In short, we extend Web Prolog into a language for
distributed programming.

The model draws on Erlang's \textit{network-transparent concurrency}: if we
know a node's identity and we are authorised, we can spawn an actor
there; if we know an actor's pid or its registered name, we can send it a message, even across node boundaries. Unlike an Erlang cluster, however, the
Prolog Web is as open as the Web itself. Here, \emph{openness} means more than
mere reachability. It means heterogeneous ownership, weak trust assumptions,
communication over ordinary Web infrastructure, and an environment in which
latency, disconnects, and partial failure are part of normal operation.
Openness therefore brings both opportunity and risk, so the chapter pays close
attention to safety: how lifetime discipline and capability separation prevent
orphan processes and resource leaks.


Openness also forces a basic separation of authority. In the Prolog
Web there are two roles: \emph{clients}, who temporarily borrow
execution resources in order to run queries and coordinate
short-lived concurrent work, and \emph{node owners}, who are
responsible for the node as an ongoing public resource. For clients,
the default discipline is \textit{containment-by-default}: spawned processes
remain tied to the client session and are torn down when the session
ends, preventing orphans and resource leaks. For owners, the
corresponding power is \textit{provisioning}: the ability to install
node-resident predicates and start long-lived actor services meant
to be shared and kept available. In practice such services are often
managed under supervision, but the generic supervision patterns that
support this belong to Chapter~\ref{sec:implementing-prolog-generics}.
This role split will recur in concrete form when we discuss remote
placement and the various code-shipping mechanisms via the
\texttt{load\_*} predicates.

The chapter is organised in five sections, moving from the concurrent
foundation through progressively higher layers of abstraction. The
concurrent Prolog Web
(Section~\ref{sec:concurrent-prolog-web}) extends the actor model
of Chapters~\ref{sec:web-prolog} and~\ref{sec:prolog-agents} across
node boundaries: remote spawning, cross-node messaging, orphan
prevention, and the unified semantics that treats local and remote
communication as instances of the same model. With the concurrent
substrate in place, we can already speak concretely about node-resident
services: long-lived actors started by the node owner, published under
stable names, and invoked by clients over the Prolog Web. The reusable
engineering patterns that structure such services -- servers,
supervisors, and other generic behaviours -- are introduced in the
next chapter and implemented in Chapter~\ref{sec:implementing-prolog-generics}.

The sequential Prolog Web
(Section~\ref{sec:sequential-prolog-web}) then introduces a
synchronous, nondeterministic abstraction built on top of the
concurrent layer. The central construct, \texttt{rpc/2-3}, can be
understood as a disciplined use of remote toplevel actors, though
in practice it is implemented over stateless HTTP backed by cached
toplevels operating in save-and-resume mode. The dependency is
genuine: the sequential interface hides the actor machinery, but it
is always, at bottom, built from it.

The programmable Prolog Web
(Section~\ref{sec:programmable-prolog-web}) shows how code and data
flow between client and node: clients may inject task-specific
predicates into their own execution contexts, while node owners
install predicates and service actors that define what the node
offers.

A discussion of timing and failure boundaries
(Section~\ref{sec:timing-failure-boundaries}) makes explicit where
waiting becomes observable and how timeouts attach to different
interaction models. Finally, a section on security
(Section~\ref{sec:security-prolog-web}) asks what makes it safe to
let strangers run code on a node, and answers in terms of three
structural mechanisms -- inexpressibility, invisibility, and
containment -- together with the deployment concerns that complete
the picture. Chapter~\ref{sec:paradigm} later returns to the
architectural and paradigmatic questions that these technical
sections raise.

%\chapter{The Prolog Web}\label{sec:the-prolog-web}
%
%\begin{quotation}
%\noindent Imagine the Web wrapped in Prolog, running on top of a distributed architecture -- a network of nodes exposing HTTP and WebSocket APIs, speaking web formats such as JSON. Think of it as a high-level Web, with Prolog agents capable of serving answers to queries -- answers that follow from what the Web knows. Moreover, imagine it being programmable, allowing Web Prolog source code to flow in either direction: from client to node, or from node to client. This is what \textbf{the Prolog Web} is all about.
%\vspace{-0.2cm}
%\flushright\textit{The Prolog Web -- the elevator pitch}
%\end{quotation}
%\vspace{0.5cm}
%
%\noindent The previous chapters introduced Web Prolog as a language
%and Prolog agents as programmable entities. This chapter turns to
%the third element of the Prolog Trinity: the \emph{Prolog Web}
%itself -- a network of nodes hosting actors that communicate across
%node boundaries. In short, we extend Web Prolog into a language for
%distributed programming.
%
%The model draws on Erlang's \textit{network-transparent concurrency}: if we
%know a node's identity and we are authorised, we can spawn an actor
%there; if we know an actor's pid or its registered name, we can send it a message, even across node boundaries. Unlike an Erlang cluster, however, the
%Prolog Web is as open as the Web itself. Openness brings both opportunity and risk, so the chapter pays close attention to safety: how lifetime discipline and capability separation prevent orphan processes and resource leaks.
%
%Openness also forces a basic separation of authority. In the Prolog
%Web there are two roles: \emph{clients}, who temporarily borrow
%execution resources in order to run queries and coordinate
%short-lived concurrent work, and \emph{node owners}, who are
%responsible for the node as an ongoing public resource. For clients,
%the default discipline is \textit{containment-by-default}: spawned processes
%remain tied to the client session and are torn down when the session
%ends, preventing orphans and resource leaks. For owners, the
%corresponding power is \textit{provisioning}: the ability to install
%node-resident predicates and start long-lived actor services meant
%to be shared and kept available, typically under supervision. This
%role split will recur in concrete form when we discuss remote
%placement and the various
%code-shipping mechanisms via the \texttt{load\_*} predicates.
%
%The chapter is organised in five sections, moving from the concurrent
%foundation through progressively higher layers of abstraction. The
%concurrent Prolog Web
%(Section~\ref{sec:concurrent-prolog-web}) extends the actor model
%of Chapters~\ref{sec:web-prolog} and~\ref{sec:prolog-agents} across
%node boundaries: remote spawning, cross-node messaging, orphan
%prevention, and the unified semantics that treats local and remote
%communication as instances of the same model. With the concurrent
%substrate in place, the section also shows how generic behaviours --
%servers and supervisors -- can be used to structure node-resident
%services; their implementations are deferred to
%Chapter~\ref{sec:implementing-prolog-generics}, but here we demonstrate their use as
%building blocks for robust distributed systems.
%
%The sequential Prolog Web
%(Section~\ref{sec:sequential-prolog-web}) then introduces a
%synchronous, nondeterministic abstraction built on top of the
%concurrent layer. The central construct, \texttt{rpc/2-3}, can be
%understood as a disciplined use of remote toplevel actors, though
%in practice it is implemented over stateless HTTP backed by cached
%toplevels operating in save-and-resume mode. The dependency is
%genuine: the sequential interface hides the actor machinery, but it
%is always, at bottom, built from it.
%
%The programmable Prolog Web
%(Section~\ref{sec:programmable-prolog-web}) shows how code and data
%flow between client and node: clients may inject task-specific
%predicates into their own execution contexts, while node owners
%install predicates and service actors that define what the node
%offers.
%
%A discussion of timing and failure boundaries
%(Section~\ref{sec:timing-failure-boundaries}) makes explicit where
%waiting becomes observable and how timeouts attach to different
%interaction models. Finally, a section on security
%(Section~\ref{sec:security-prolog-web}) asks what makes it safe to
%let strangers run code on a node, and answers in terms of three
%structural mechanisms -- inexpressibility, invisibility, and
%containment -- together with the deployment concerns that complete
%the picture. Chapter~\ref{sec:paradigm} later returns to the
%architectural and paradigmatic questions that these technical
%sections raise.


\section{The concurrent Prolog Web}\label{sec:concurrent-prolog-web}

As introduced above, all the mechanisms described in the preceding
chapters -- spawning, sending, linking, monitoring, registering --
work transparently across node boundaries. This section makes that
claim concrete by introducing remote placement via the \texttt{node}
option and showing how actor-style interaction spans nodes. The focus
is on \emph{where} an actor runs and how it exchanges messages across
the network, but also on the safety and authority questions that
network-transparent concurrency raises on an open Web.
Programmability mechanisms such as \texttt{load\_*} options are deliberately postponed to
Section~\ref{sec:programmable-prolog-web}.



\subsection{Actor talking to remote actor}\label{sec:talking-to-remote-actor}

For specifying the node on which we wish to create an actor process, \texttt{spawn/3} can be passed the \texttt{node} option: 

\begin{verbatim}
?- spawn(echo_actor, Pid, [
       node('http://n8.org'),
       monitor(true)
   ]),
   register(echo_actor, Pid).
Pid = 99231321@'http://n8.org'.
?- 
\end{verbatim}

\noindent Note that the pid is a compound term: an integer paired with a URI using the infix operator \texttt{@}. If \texttt{n8.org} is a node conforming to the ACTOR profile, and if the node's shared database defines the predicate \texttt{echo\_actor/0}, then we can have a conversation with our remote echo server like so:

\begin{verbatim}
?- self(Self),
   echo_actor ! echo(Self, hello).
Self = 71773120.
?- flush.
Shell got echo(hello)
true.
?- whereis(echo_actor, Pid),
   exit(Pid, die).
Pid = 99231321@'http://n8.org'.
?- receive({A -> true}).
A = down(99231321@'http://n8.org',99231321@'http://n8.org',die).
?- 
\end{verbatim}


\subsection{Actors playing ping-pong while on different nodes}\label{sec:pong-pong-two-nodes}

The actors playing ping-pong back in Chapter~\ref{sec:web-prolog} kept their game inside one node. Passing the option \texttt{node('http://n2.org')} to one of the calls to \texttt{spawn/3} allows the game to be played between two nodes:

\begin{verbatim}
ping_pong :-
    spawn(pong, Pong_Pid, [
        node('http://n2.org')
    ]),
    spawn(ping(3, Pong_Pid)).
\end{verbatim}

\noindent Figure~\ref{fig:ping-pong-scenario} illustrates this scenario.

\begin{figure}[h]
    \centering
	\includegraphics[width=\textwidth]{ping-pong-scenario}
    \caption{Actors playing ping-pong while on different nodes.}
    \label{fig:ping-pong-scenario}
\end{figure}

\noindent But why do we get the output written by the ``pinger'' only, and not the output from the ``ponger''? The asymmetry in the output is not an accident of the example but a consequence of a deliberate capability rule. A terminal is attached to one particular toplevel actor, and rendering output in that terminal is treated as an \emph{I/O capability} that is scoped to the toplevel's local process lineage. Concretely, only actors that (i) run on the same node as the toplevel to which the terminal is attached, and (ii) are local descendants of that toplevel, may write output that is forwarded to the client for display. In any case, readers may rest assured that the actors really \textit{do} play ping-pong.


\subsection{The node option works for toplevels too!}\label{sec:node-option-for-toplevels}

Toplevels are actors, so it should come as no surprise that the \texttt{node} option can also be used to create a toplevel process on a remote node: 

\begin{verbatim}
?- toplevel_spawn(Pid, [
       node('http://n1.org')
   ]).
Pid = 83110912@'http://n1.org'.
?-   
\end{verbatim}

\noindent The query \texttt{?-human(Who)} executes on the remote node \texttt{n1.org}, with answers returned to the local shell:

\begin{verbatim}
?- toplevel_call($Pid, human(Who), [
       template(Who),
       limit(1)
   ]).
true.
?- receive({Msg -> true}).
Msg = success(83110912@'http://n1.org',[plato],true).
?- toplevel_next($Pid).
true.
?- receive({Msg -> true}).
Msg = success(83110912@'http://n1.org',[aristotle],false).
?-
\end{verbatim}

\noindent Network transparency makes distributed concurrency easy to \emph{use}. The harder question is how to make it \textit{safe} on an open, failure-prone Web: how do we prevent orphan processes, constrain resource usage, and ensure that failure leads to reclamation rather than silent accumulation? The answer is largely about lifetime discipline, supervision, and the separation of authority between clients and node owners.

\subsection{Node-resident actor processes}\label{sec:node-resident-actors}

In addition to predicates in the shared database, the owner of a
node may install \textit{node-resident actor processes} -- long-lived
actors that are intended to be durable, shared, and continuously
available. They are usually of a kind that we do not want
unauthorized external clients to be able to create or destroy; doing
this is the privilege of the node owner. This is also where supervision trees become most useful: crash recovery for node-resident services is part of the node owner's responsibility, and Chapter~\ref{sec:implementing-prolog-generics} introduces the supervisor behaviour that makes this responsibility tractable.


\subsubsection{A node-resident count actor}

We begin with a deliberately small example. Assume that the owner of the node \texttt{n1.org} starts a count actor and gives it a stable registered name:

\begin{verbatim}
?- spawn(count_actor(0), Pid),
   register(counter, Pid).
\end{verbatim}

\noindent The registered name can now be used instead of the pid when sending messages to the process:

\begin{verbatim}
?- self(Self).
Self = 51230945.
?- counter@'http://n1.org' ! count($Self),
   receive({Count -> true}).
Count = 1.
?- counter@'http://n1.org' ! count($Self),
   receive({Count -> true}).
Count = 3.
?-
\end{verbatim}

\noindent The example is intentionally simple, but it already shows an important shift in perspective. This actor is not a client-local helper process created for one temporary task; it is a node-resident process with a stable public name, and it therefore behaves like a shared stateful resource. The fact that the second reply is \texttt{3} rather than \texttt{2} is not an accident of presentation, but evidence that the actor's state is shared across interactions: another client has incremented the counter in the meantime.

In other words, once a node owner chooses to run and name an
actor in this way, the actor begins to look less like a private process
and more like a service. The next subsection develops this idea further
by showing a publish-subscribe actor whose service character is even more
explicit.

\subsubsection{A node-resident publish-subscribe service}

The same service perspective becomes even clearer in a publish-subscribe
example. Consider the following actor:

\begin{verbatim}
pubsub_actor(Subscribers0) :-
    receive({
        publish(Msg) ->
            forall(member(Pid, Subscribers0), Pid ! msg(Msg)),
            pubsub_actor(Subscribers0);
        subscribe(Pid) ->
            pubsub_actor([Pid|Subscribers0]);
        unsubscribe(Pid) ->
            (   select(Pid, Subscribers0, Subscribers)
            ->  pubsub_actor(Subscribers)
            ;   pubsub_actor(Subscribers0)
            )
    }).
\end{verbatim}

\noindent This actor maintains a list of subscribers. A
\texttt{subscribe(Pid)} message adds a subscriber, an
\texttt{unsubscribe(Pid)} message removes one, and a
\texttt{publish(Msg)} message broadcasts \texttt{Msg} to all
currently subscribed processes. On its own, this is just another
actor definition. It becomes a \emph{service} when the node owner
chooses to start it as a durable node-resident process and publish it
under a stable name:

\begin{verbatim}
?- spawn(pubsub_actor([]), Pid),
   register(pubsub_service, Pid).
\end{verbatim}

\noindent A client can then subscribe to the service and wait for
messages from it:

\begin{verbatim}
?- self(Self),
   pubsub_service@'http://n8.org' ! subscribe(Self),
   repeat,
   writeln("Waiting for a message ..."),
   receive({
       msg(Message) ->
           format("Received: ~p~n", [Message]),
           fail
   }).
Waiting for a message ...
Received: hello
Waiting for a message ...
\end{verbatim}

\noindent The message \texttt{hello} was received because some client
published it to the same named actor:

\begin{verbatim}
?- pubsub_service@'http://n8.org' ! publish(hello).
true.
?-
\end{verbatim}

\noindent The important point is not the specific protocol, but the
change in status. Once the actor is started by the node owner,
registered under a stable name, and intended for repeated use by
multiple independent clients, it should be understood as part of the
node's public surface. It is no longer merely a process that happens
to be running; it is a shared conversational resource offered by the
node.

This example also illustrates a useful contrast with the
counter actor. The counter service exposes shared \emph{state}: what
one client observes depends on what other clients have already done.
The publish-subscribe service instead exposes shared
\emph{coordination}: clients interact indirectly through a common
mediating process. In both cases, however, the essential pattern is
the same. A node-resident actor with a stable public name behaves as
a service.

The implementation shown here is intentionally minimal. In
practice, one would typically place such a service under supervision,
possibly persist subscriptions in a more durable store, and perhaps
attach access control or topic structure. Those are engineering
refinements. The conceptual point is already visible in the small
example: named node-resident actors provide a natural way to realise
stateful services on the Prolog Web.

\subsubsection{Naming, discovery, and service publication}
\label{sec:naming-discovery}

The preceding examples rely on stable service names -- \texttt{counter},
\texttt{pubsub\_service} -- to make node-resident actors reachable by
clients. In this design, these are not ordinary client-managed
registrations. A client may still use \texttt{register/2},
\texttt{whereis/2}, and \texttt{unregister/1} for actors that it has
itself created and named, but published services live in a separate
owner-controlled service registry. Clients may send to a published
service name, but they do not thereby gain the ability to discover the
service's pid or to remove its published name. This raises the practical
question: how do clients discover which services a node makes available?

The standard design principle is to separate a stable,
location-independent identifier from the service's current location
\cite{vanSteenTanenbaum2023DistributedSystems}. The Prolog Web inherits
this directly: a service name should not bind clients to a particular
pid, host, transport endpoint, or deployment topology. Services may be
restarted, moved, replicated, or versioned; the published name remains
stable.

This also clarifies the distinction between ordinary naming and service
publication. Ordinary registration is a convenience for private or
client-local actors. Service publication is a different act: the node
owner places a name into the node's public service namespace. An
implementation should therefore provide dedicated operations --
\texttt{register\_service/2} and \texttt{unregister\_service/1} -- for
published services, leaving the ordinary naming predicates available for
client-managed actors.

The safe pattern is to treat discovery as an owner-curated
\emph{publication mechanism}: the node exposes a directory that lists
only those services intended as part of its public surface. Discovery
should resolve \emph{published} names, not reveal the existence of
everything that happens to be running.

This is a design invariant, not merely a preference. If clients could
enumerate all running processes, or if ordinary client naming operations
exposed or removed published services, discovery would collapse into
introspection and name discipline would be lost. The security
implications of this principle are developed in
Section~\ref{sec:invisibility}.

In the Trinity, the cleanest way to expose such a directory is to make it
queryable through \texttt{rpc/2-3}. Concretely, a node that wishes to
advertise services may provide a relation such as \texttt{service/2} (or
\texttt{service/3}) in its shared database, for example:

\begin{verbatim}
service(counter, meta(actor, protocol(count_v1))).
service(pubsub_service, meta(actor, protocol(pubsub_v1))).
\end{verbatim}

\noindent A client can then discover advertised services using an
ordinary goal-oriented query:

\begin{verbatim}
?- rpc('http://n1.org', service(Name, Meta)).
Name = counter,
Meta = meta(actor, protocol(count_v1)) ;
Name = pubsub_service,
Meta = meta(actor, protocol(pubsub_v1)).
?-
\end{verbatim}

\noindent This keeps the design uniform. Service discovery remains
declarative and compositional: clients can ask for all services, filter
by protocol, restrict to public entries, or combine discovery with other
constraints in the same query. At the same time, it does not weaken the
security model, because the node remains free to decide what it
advertises, to require authorisation for the directory predicate, and to
ensure that the directory reveals only what it has chosen to publish.

\medskip
\noindent With this distinction in hand, the earlier counter example can
be read more precisely. Once an actor is published under a stable service
name such as \texttt{counter}, it is no longer merely a running process
known to its owner, but part of the node's public surface. In that role
it behaves as a shared service: multiple clients may interact with it
through the same published name, while the underlying pid remains an
implementation detail of the node. Observations such as receiving
\texttt{3} rather than \texttt{2} are then best understood not as
oddities, but as visible consequences of shared state.



\subsection{Uniform messaging across node boundaries}\label{sec:uniform-messaging}

The preceding examples show that the same primitives -- spawning,
sending, linking, monitoring, registering -- work whether the target
process is local or remote.  This uniformity is not accidental; it
rests on a deliberate design choice inherited from the unified
semantics proposed by \citet{conf/erlang/SvenssonFE10} for
Erlang-like systems.  The key idea is to eliminate the semantic
distinction between local and remote communication: all
side-effecting operations are mediated by a node controller
abstraction, so that placement becomes a configuration detail rather
than a semantic one. Chapter~\ref{sec:paradigm} returns to the deeper implications of this distribution-first stance; here we spell out the messaging
model that makes it concrete.


\subsubsection{Two-step messaging model}
Following Svensson et al.~\cite{conf/erlang/SvenssonFE10}'s \emph{ether}
abstraction, we model messaging as a two-step mechanism:
\emph{sending} places an in-flight record into a virtual system
queue (the \emph{ether}), and a later \emph{delivery} step appends
the message to the receiver's mailbox. This separation allows us to
define program meaning independently of transport.

The in-flight layer is a semantic device rather than an
implementation requirement. Locally, a runtime may realize sending
by inserting the message directly into the receiver's mailbox, which
can be seen as choosing a schedule in which each send is immediately
followed by its delivery. In Web Prolog, the in-flight phase may be
implemented using multiple transports, as illustrated below.

\begin{center}
\begin{tabular}{c}
Process $A$ on node $N_1$ \\
$\downarrow$ \quad \texttt{B ! M} \\
\textbf{in-flight record} $(A, B, M)$ in the ether \\
$\downarrow$ \quad (transport: local queue \;|\; WebSocket \;|\; HTTP \;|\; relay) \\
Receiver mailbox of $B$ on node $N_k$ \\
$\downarrow$ \quad \texttt{receive(\{...\})}
\end{tabular}
\end{center}

\noindent The transport realizations in Web Prolog are as follows.

\begin{itemize}
\item \textbf{Local delivery}: the runtime enqueues $(A,B,M)$ in an
internal scheduler structure and later appends $M$ to $B$'s mailbox.
\item \textbf{Direct remote delivery}: $(A,B,M)$ is serialized and
shipped from $N_1$ to $N_k$ as a WebSocket frame and then enqueued
for mailbox delivery.
\item \textbf{Request-based delivery}: $(A,B,M)$ is carried in an
HTTP request to $N_k$ (or to a node gateway), which authenticates,
deserializes, and enqueues the message.
\item \textbf{Relay-mediated delivery}: $(A,B,M)$ is sent to an
intermediary that forwards it to $N_k$, after which normal mailbox
delivery occurs.
\end{itemize}

\noindent This separation lets us specify one uniform concurrency
model while allowing multiple deployment choices. Ordering guarantees
are stated over \emph{delivery} into mailboxes: per sender-receiver
FIFO when sender and receiver are addressed without mixing addressing
modes, with no global order across different senders
(Section~\ref{sec:message-passing-guarantees}). Liveness properties
depend on the fairness of the delivery schedule rather than on the
specifics of the network path.



\section{The sequential Prolog Web}\label{sec:sequential-prolog-web}

The \textit{sequential Prolog Web} is the fragment of the Prolog Web
in which every client interaction can be explained by ordinary
Prolog-style backtracking. More precisely, a node belongs to the
sequential Prolog Web if, for every client session, the sequence of
answer substitutions returned for each goal is exactly the sequence
that a single-threaded SLD-style execution would produce. No
reasoning about interleavings, mailboxes, or concurrent actors is
needed to understand what happens -- the entire observable behaviour
admits explanation by a single resumable logical history.

The \textit{concurrent} partition is fundamentally different: it
admits multiple simultaneous causal histories, requires reasoning
about interleavings, and cannot in general be reduced to a single
sequential trace. Yet the two partitions are not independent. The
sequential layer of the Prolog Web is built directly on top of the concurrent one, and
understanding that dependency is the key to understanding why the
sequential interface looks and behaves as it does.

The full layering, from bottom to top, is as follows:

\begin{enumerate}
    \item Actor processes and asynchronous messagin.
    \item Toplevel actors, which add goal evaluation and answer-stream management to the basic actor model.
    \item \texttt{rpc/2-3} over a remote toplevel -- the conceptual implementation that makes the synchronous, nondeterministic interface explicit.
    \item \texttt{rpc/2-3} over stateless HTTP -- the practical implementation, backed by cached toplevels operating in save-and-resume mode, which realizes the same interface more efficiently.
\end{enumerate}

\noindent Items 3 and 4 expose the same interface; they differ in realization, not in semantics. Each step in the stack thins the interface and hides more of the underlying machinery.

Because the concurrent layer has already been established, we can now present \texttt{rpc/2-3} and its companions \texttt{promise/3-4} and \texttt{yield/2-3} with the reader already understanding what lies beneath. Not every distributed program needs explicit concurrency -- often a caller simply wants to ask another node to solve a goal and wait for answers -- but the sequential interface is always, at bottom, a disciplined use of the concurrent one.


\subsection{Nondeterministic remote procedure calls}\label{sec:non-deterministic-procedure-calls}

In Web Prolog, \texttt{rpc/2-3} is a high-level meta-predicate for making \emph{nondeterministic remote procedure calls}. It allows a process running on a node $N_1$ to call and attempt to solve a goal in the Prolog context of another node $N_2$, taking advantage of the programs and data offered by $N_2$ as if they were local to $N_1$. Such calls are synchronous: the caller blocks while waiting for answers. In return, \texttt{rpc/2-3} is remarkably easy to use for distributing programs across two or more nodes. It can therefore be seen as a construct for distributed programming that trades concurrency for ease of programming.

A Web Prolog client performs a remote call by providing a URI identifying the node and a goal to be solved in that node's context. The optional third argument is a list of options. Here is a simple example:

\begin{verbatim}
?- rpc('http://n1.org', human(Who)).
Who = plato ;
Who = aristotle.
?-
\end{verbatim}

\noindent Note that solutions are returned on demand and in the familiar one-solution-at-a-time fashion: variable bindings appear as instantiations of variables occurring in the goal.

\subsubsection{Using rpc/3 with the limit option}\label{sec:rpc-with-limit}

The \texttt{rpc/3} predicate supports a \texttt{limit} option, inherited from
the underlying \texttt{toplevel\_call/3} predicate. Passing \texttt{limit(K)}
instructs the remote node to compute at most $K$ solutions per network roundtrip.
The default is \texttt{infinite}, meaning all solutions are fetched in a single
exchange. Here is what happens if we override that default with \texttt{limit(1)}:
\begin{verbatim}
?- rpc('http://n1.org', human(Who), [
       limit(1)
   ]).
Who = plato ;
Who = aristotle.
?-
\end{verbatim}
\noindent From the caller's perspective, the result is identical. What changed is
invisible: instead of retrieving both solutions in a single roundtrip, the call
now required \emph{two} -- one per solution. In this case, the option adds cost
without benefit. Later, we shall see contexts in which limiting the number of
solutions per roundtrip matters in a genuinely positive way.

\medskip
\noindent In addition to the \texttt{limit} option, \texttt{rpc/3} inherits the \texttt{load\_*} options from \texttt{spawn/3} and the \texttt{timeout} and \texttt{once} options from the toplevel mechanism. Examples using \texttt{load\_*} options will be given in Section~\ref{sec:injecting-code}. See Appendix~\ref{sec:manual} for details about the other options.


\subsubsection{The semantics of rpc/2-3}\label{sec:semantics-rpc}

For intuition it is useful to note that a call to
\texttt{rpc(URI,\,Goal)} is semantically equivalent to the following
pattern, which makes explicit that solutions arrive as instantiations
of a copy that is then unified back into the original goal:

\begin{verbatim}
copy_term(Goal, Copy),
call(Copy),           % executed on node at URI
Goal = Copy.
\end{verbatim}

\noindent The key implication is that no variables are shared across
nodes: the remote execution operates on a renamed-apart copy, and
only the induced substitution is returned to the caller via
unification with the original goal.

This three-line characterisation is intended as a \emph{logical}
specification of \texttt{rpc/2-3}: it defines what answers the
caller may observe, not how those answers are obtained. As a specification it is precise and complete for the intended domain. Web Prolog restricts remote goals to pure Herbrand terms -- attributed variables are not supported, and the serialisation protocol ensures that bindings returned to the caller are always meaningful in the caller's context, ruling out node-local references such as stream handles or database pointers. What the characterisation deliberately abstracts over is the operational machinery required to deliver those answers across a network: the creation of a remote process, the protocol that sustains nondeterminism across multiple solutions, and the handling of timeouts, exceptions, and partial failure. These concerns are addressed by actor infrastructure described in the sections that follow.


\medskip
\noindent A notable property of \texttt{rpc/2-3} is that it preserves the
logical purity of the goal it calls: if the goal is pure, then the
answers returned by \texttt{rpc/2-3} are exactly the logical
consequences of the remote database -- the network introduces no
additional bindings, side effects, or observable state
changes.\footnote{This is a property it shares with the
\texttt{call/N} family of Prolog predicates.} This holds at the
level of answer substitutions, assuming successful communication; in
practice, timeouts, connection failures, and resource limits may
prevent all answers from being delivered, but they do not corrupt the
answers that are delivered. Later, we return to this observation when
discussing the idea of a \emph{pure} Prolog Web.

\subsection{Implementing \texttt{rpc/2-3} on top of a toplevel actor}\label{sec:rpc-on-top-of-toplevel}

The central construct of the sequential layer, \texttt{rpc/2-3}, can be
implemented straightforwardly on top of a toplevel actor: spawn a toplevel on
the target node, submit a goal, and collect answer messages in a loop. From the
caller's perspective the interaction is synchronous and nondeterministic --
variables are bound on success, backtracking yields the next solution -- but
underneath, the familiar PTCP communication protocol is doing the work:

\begin{verbatim}
rpc(URI, Goal) :-
    rpc(URI, Goal, []).
    
rpc(URI, Goal, Options) :-
    toplevel_spawn(Pid, [
         node(URI),
         session(false),
         monitor(false),
       | Options
    ]),
    toplevel_call(Pid, Goal, Options),
    wait_answer(Pid, Goal).
\end{verbatim}

\noindent The remote toplevel is spawned with \texttt{session(false)} so that it
terminates after the query has been exhausted. The implementation does not monitor
the spawned process, so no \texttt{down} message needs to be handled. The loop
can therefore terminate cleanly on failure, error, or after the last slice has
been consumed:

\begin{verbatim}
wait_answer(Pid, Goal) :-
    receive({
        success(Pid, Slice, true) ->
            (   member(Goal, Slice)
            ;   toplevel_next(Pid),
                wait_answer(Pid, Goal)
            ) ;
        success(Pid, Slice, false) ->
            member(Goal, Slice) ;
        failure(Pid) -> !, fail ;
        error(Pid, Error) ->
            throw(Error)
    }).
\end{verbatim}

\noindent The nondeterministic behaviour arises from two sources. An answer term of the form \texttt{success(Pid,\,Slice,\,\ldots)} message contains a \emph{slice} of
solutions, and \texttt{member(Goal,\,Slice)} enumerates them without further
network communication. Only when a slice is exhausted does execution fall through
to the second disjunct, which requests a new slice via \texttt{toplevel\_next/1}.
In this way, deterministic message exchanges with a remote toplevel are turned
into a nondeterministic, on-demand interface at the caller.

Since the caller and callee do not perform independent work during overlapping
periods, this is a purely synchronous and sequential pattern: \texttt{rpc/2-3}
contributes to the sequential fragment of the Prolog Web.





\subsection{Implementing \texttt{rpc/2-3} on top of the stateless HTTP API}
\label{sec:rpc-on-http}

Looking at the implementation of \texttt{rpc/2-3} in terms of the lower-level
\texttt{toplevel\_*} predicates makes visible what the high-level specification
deliberately abstracts over. In particular, it becomes clear that all
communication follows a strict request-response alternation: the caller creates
a query remotely, then repeatedly asks for the next answer, blocking each time
until a response arrives. There are no unsolicited messages, no server-initiated
notifications, no concurrent streams -- just a disciplined exchange in which
every message from the client is answered by exactly one message from the server.
This is precisely the traffic pattern that a stateless HTTP API is designed to
carry, and it is why \texttt{rpc/2-3} can be implemented entirely on top of the
sequential layer's HTTP interface without any loss of expressiveness.

This implementation is not merely adequate but arguably preferable. Because each
exchange is a self-contained HTTP request-response pair, there is less risk of
leaving orphan queries behind on the remote node. The approach works between any
two ISOBASE-conforming nodes regardless of whether they support the programmable
layer. It passes cleanly through firewalls, proxies, and load balancers. Failures
are unambiguous -- every request either receives a response or times out. And the
entire conversation is inspectable with standard HTTP tooling. In short, for the
synchronous, one-answer-at-a-time pattern that \texttt{rpc/2-3} requires,
stateless HTTP is not a compromise but the right fit.

Compared to the WebSocket-based implementation shown above, the HTTP version is
more involved. The additional complexity lies not in the control structure, which
is shared between the two, but in transport concerns: serialising goals and
templates, encoding terms correctly, and handling authentication, redirects, TLS
settings, and timeouts. The full SWI-Prolog implementation, including URL
construction, term encoding, and HTTP option handling, is given in Appendix~B.



\subsection{Browser talking to a node talking to a node}\label{sec:example-two-hop}

\noindent In the scenario depicted in Figure~\ref{fig:interaction-isobase}, a user of a Prolog terminal running in a browser interacts with a remote node on the Prolog Web, which in turn interacts with another node. 

\begin{figure*}[h!]
    \centering
    \includegraphics[width=\textwidth]{interaction-isobase.png}
    \caption{Stateless interaction in a scenario involving a web browser and two remote nodes.}
    \label{fig:interaction-isobase}
\end{figure*}

The solutions shown in the terminal follow from the \emph{combination} of two programs: the code on the upper left in Figure~\ref{fig:interaction-isobase} belongs to \texttt{n2.org}, and it contains a URI pointing to the code offered by \texttt{n1.org}. The browser-to-node interaction as well as the node-to-node interaction take place over HTTP and are stateless in the sense that each request can be understood and processed without relying on server-side session state.


To retrieve the first solution to \texttt{?-mortal(Who)}, the browser issues a request to \texttt{n2.org} asking for exactly one solution:

\begin{verbatim}
GET http://n2.org/call?goal=mortal(Who)&limit=1
\end{verbatim}

\noindent Since the user requests solutions one-by-one, the response contains a single binding together with a flag indicating that more solutions exist:

\begin{verbatim}
{ "type": "success",
  "data": [{"Who":"socrates"}],
  "more": true }
\end{verbatim}



\noindent When the user types a semicolon, the browser requests the next solution. In this example, producing that next solution requires \texttt{n2.org} to consult \texttt{n1.org}. Concretely, \texttt{n2.org} evaluates \texttt{mortal/1} and, as part of doing so, triggers an \texttt{rpc/2-3} call for \texttt{human(Who)} against \texttt{n1.org}. Operationally, this involves an HTTP request to \texttt{n1.org} such as:

\begin{verbatim}
GET http://n1.org/call?goal=human(Who)&format=prolog
\end{verbatim}

\noindent Because of the default values of the corresponding slicing parameters, the response in this case contains a complete slice with all solutions to \texttt{human(Who)} currently returned in one roundtrip. Since the URI specified \texttt{format\,=\,prolog}, this was the result:

\begin{verbatim}
success([human(plato),human(aristotle)], false)
\end{verbatim}

\noindent At this point the implementation of \texttt{rpc/2-3} on \texttt{n2.org} must convert the returned slice into the intended nondeterministic behaviour. The usual mechanism is to keep the slice locally and enumerate it using \texttt{member/2}. This yields solutions to \texttt{human(Who)} on backtracking without further network communication, until the slice is exhausted.

Since the browser client still wants only one solution at a time, \texttt{n2.org} returns only the first derived solution, \texttt{Who=plato}, to the browser:

\begin{verbatim}
{ "type": "success",
  "data": [{"Who":"plato"}],
  "more": true }
\end{verbatim}


\noindent When the user again types a semicolon, the browser requests another solution to \texttt{mortal(Who)}. In this particular run, no additional request to \texttt{n1.org} is needed: \texttt{n2.org} can simply backtrack within the local \texttt{member/2} enumeration of the slice already obtained from \texttt{n1.org}. The browser request therefore looks like:

\begin{verbatim}
GET http://n2.org/call?goal=mortal(Who)&offset=2&limit=1
\end{verbatim}

\noindent and the response contains the next binding, \texttt{Who=aristotle}, now with \texttt{more:false}:

\begin{verbatim}
{ "type": "success",
  "data": [{"Who":"aristotle"}],
  "more": false }
\end{verbatim}

\noindent This completes the description of what happens ``under the hood'' in the scenario of Figure~\ref{fig:interaction-isobase}: the browser drives one-solution-at-a-time enumeration through stateless HTTP requests to \texttt{n2.org}, while \texttt{n2.org} uses \texttt{rpc/2-3} to consult \texttt{n1.org} and reconstructs nondeterminism locally from slices of answers.



\subsubsection{Large result sets and the \texttt{limit} option}

If \texttt{n1.org} had millions of clauses of \texttt{human/1},
returning a huge slice in one response would be undesirable. In that
case \texttt{n2.org} can pass a smaller \texttt{limit} to
\texttt{rpc/2-3} to retrieve answers in manageable batches:

\begin{verbatim}
human(Who) :-
    rpc('http://n1.org', human(Who), [
        limit(1000)
    ]).
\end{verbatim}

\noindent This avoids swamping the caller with data. A goal with
very many solutions will be delivered incrementally over multiple
roundtrips, but since \texttt{rpc/2-3} is nondeterministic, the call
remains logically complete provided the caller keeps asking for more.


\medskip
\noindent As this scenario involves only two nodes, the Prolog Web here is tiny. Later we argue that the same principles can, in principle, scale to arbitrarily many nodes. Key is the combination of statelessness (for robustness and simplicity) with an execution model that avoids shipping excessive data and avoids recomputing answers unnecessarily.


\subsection{Backtracking as pagination}
\label{sec:backtracking-as-pagination}

The \texttt{offset} and \texttt{limit} parameters in the stateless
HTTP API may look like a polite concession to web conventions -- a
way to make Prolog's multiple solutions fit a request--response
protocol. But the connection is deeper than that. Prolog
backtracking and HTTP pagination instantiate the same abstract idea:
incremental enumeration under external control.

A paginated endpoint does not ``return the result set''; it exposes
a procedure for generating results and a protocol for consuming them
in slices. The server produces a page, then stops. A later request
asks for the next page, which amounts to resuming enumeration at a
recorded position. Prolog behaves in the same general way. A goal
may have many solutions; the toplevel presents the next one (or the
next batch) and then waits. Internally, evaluation proceeds by
creating choice points, and enumeration can be resumed from a
well-defined position. In an interactive toplevel, that position is
represented by retained search state, and the user's semicolon is
simply the ``resume'' signal.


\subsubsection{Offset surface, cursor semantics}

Web API pagination comes in two principal forms. In \emph{offset
pagination}, the client says ``skip~$N$, return the next~$k$''; the
server must, at least logically, regenerate the first~$N$ results
and discard them. In \emph{cursor pagination}, the server hands the
client an opaque token that encodes where to resume; the client
returns the token to request the next page, and the server picks up
where it left off. The token is meaningless to the client -- it is a
handle to internal state, not a numeric position in a pre-existing
result set.

Prolog backtracking is structurally cursor-based. A choice point is
opaque retained state that captures exactly where to resume the
search: variable bindings, unexplored branches, the position in the
clause index. The semicolon hands it back and says ``continue.'' No
re-enumeration from the top occurs; no solutions are skipped and
discarded.

On the wire, however, the stateless HTTP API speaks \texttt{offset}
and \texttt{limit} because that is idiomatic REST:
%
\begin{verbatim}
GET /call?goal=hotel(Name,paris,Rating)&offset=20&limit=10
\end{verbatim}
%
This returns solutions 21--30. But the implementation need not treat
the offset as an instruction to re-enumerate from the beginning and
discard. A node can cache a suspended toplevel process keyed by goal
and offset, so that a request for page~$k{+}1$ resumes the process
from its suspended choice point rather than replaying the search.
The offset then functions not as a distance from the start, but as a
lookup key into retained state -- in effect, a cursor in offset
clothing.


\subsubsection{Consistency and the update view}

The distinction between true offset re-execution and cached-process
resumption is not merely one of cost; it has semantic consequences.
True offset pagination -- re-executing the goal fresh for each page
-- gives the client a sliding window over whatever the database
contains at the moment of the request. If facts are asserted or
retracted between pages, solutions may be duplicated, skipped, or
reordered. Cursor-based resumption, by contrast, continues an
in-flight enumeration whose view of the database is governed by the
Prolog system's update semantics.

In the Prolog literature, this is the distinction between the
\emph{logical update view} and the \emph{immediate update view}.
Under the logical update view -- which the ISO Prolog standard mandates -- a goal sees the set of clauses that existed when it was called, regardless of later modifications. A cached toplevel process
therefore provides snapshot-consistent pagination: each page
continues a coherent enumeration even as the shared database evolves
underneath. Under the immediate update view, the running goal would
see mid-enumeration changes, weakening consistency but still
preserving the search-tree position.

The choice of update view thus determines the isolation level of
paginated enumeration, connecting a seemingly low-level
implementation decision to the same consistency questions that web
developers face when paginating over mutable data -- and that
database systems address with server-side cursors and versioned
snapshots. A Prolog system that provides the logical update view
delivers, essentially for free, the kind of stable pagination that
web backends often must engineer with care.


\subsubsection{Three points on a spectrum}

The relationship between Prolog's backtracking mechanism and HTTP
pagination can be understood as a spectrum of three API styles, each
exposing more of the underlying process model.

\begin{enumerate}
\item \textbf{Stateless offset, naive implementation.} Each request
  re-executes the goal against the current database, skips~$N$
  solutions, and returns the next batch. This is true offset
  pagination: cheap on server memory, idempotent by construction, but
  vulnerable to inconsistency if the database changes between pages.

\item \textbf{Stateless offset, cached processes.} The wire protocol
  still says \texttt{offset}/\texttt{limit}, but the server caches a
  suspended toplevel process keyed by goal and offset. Offset on the
  surface, cursor under the hood. Under the logical update view, this
  gives snapshot-consistent pagination while retaining a conventional
  REST interface.

\item \textbf{Semi-stateful API with process identifiers.} The
  server hands back an explicit process identifier, and the client
  uses it to request continuation. There is no offset at all. This is
  cursor pagination without disguise, and the closest structural
  match to interactive backtracking: the pid~\emph{is} the
  choice-point handle, and ``next'' \emph{is} the semicolon.
\end{enumerate}

The three styles differ in what the client's handle refers to.
In~(1), there is no handle; each request reconstructs everything
from scratch. In~(2), the offset is an implicit handle to a
suspended search tree -- choice points and bindings against the
shared database. In~(3), the pid is an explicit handle to a full
process: search tree, variable bindings, \emph{and the process's
private database}, including any facts it has asserted locally or
code loaded via \texttt{load\_*} options. Successive ``next''
requests do not just advance enumeration; they resume a process that
may have asserted local facts, updated counters, or built up
intermediate structures as part of its computation. This is closer to
a database session with a temporary table than to a mere cursor.


\subsubsection{When the stateless API grows stateful}

The clean separation between stateless offset pagination and stateful
pid-based resumption blurs once the stateless API accepts
\texttt{load\_*} options. A request such as
%
\begin{verbatim}
GET /call?goal=solve(X)&load_module=planner&offset=0&limit=5
\end{verbatim}
%
causes the node to spin up a process that loads the specified module
into its private database before executing the goal. That process now
carries local state -- the loaded clauses -- not just a position in a
search tree over the shared database.

If the node caches this process for subsequent page requests, the
cached handle points to a process with a populated private database:
essentially the same kind of entity that a pid refers to in the
semi-stateful API. The only difference is how the client addresses it
-- an offset key derived from the goal and its options, versus an
explicit pid.

If the node does \emph{not} cache -- naive re-execution -- then each
page request re-loads the code, re-executes from scratch, and skips.
This is not merely inefficient for enumeration; it re-runs the entire
setup for every page. And if the loaded source has changed between
requests, different pages may be enumerating over different programs
entirely.

The conclusion is that even the API that looks most RESTful is, in
the presence of \texttt{load\_*} options, implicitly creating
processes with private state. The question is only whether that state
is retained across pages (cached) or reconstructed (naive), and
whether the client gets an opaque handle to it (pid) or addresses it
indirectly (goal + options + offset as a composite key). The
semi-stateful API simply makes this honest.


\subsubsection{Caching and the web}

The stateless design has a further consequence that is easy to
overlook. Because Web Prolog pagination uses standard HTTP mechanics,
it inherits the Web's caching infrastructure. A request for
\texttt{?goal=mortal(X)\&offset=0\&limit=10} is a plain URI with a
deterministic response. CDNs, reverse proxies, and browser caches
can store and serve it without modification. Load balancers can route
requests to any backend since there is no session state to preserve.
And retry semantics are simple: the request is idempotent, so
failures can be retried safely. None of this requires any special
support from Web Prolog; it falls out automatically from the decision
to expose backtracking as stateless pagination over HTTP.

These caching benefits apply most directly to the stateless API
without \texttt{load\_*} options -- precisely the case where the
offset is closest to a true numeric skip parameter and the response
is most straightforwardly deterministic. As we move along the
spectrum toward cached processes and \texttt{load\_*} options, the
caching story becomes more nuanced, but the surface-level HTTP
compatibility is preserved.


\subsection{Programming with \texttt{rpc/2-3}}\label{sec:programming-with-rpc}

Because \texttt{rpc/2-3} is a simple, synchronous, nondeterministic interface, it is easy to build small programming schemes on top of it. We sketch a few patterns here.


\subsubsection{The parallel/1 predicate}\label{sec:parallel-pred}

A call to \texttt{parallel(Goals)} should

\begin{enumerate}
\item block until all work has been done, and no longer than necessary,
\item succeed, with variable bindings, if all goals succeed,
\item fail as quickly as possible if any goal fails, and 
\item rethrow any errors thrown by a goal, also as quickly as possible.
\end{enumerate}

\noindent In other words, running a list of goals in parallel should behave in the same way as when running them sequentially, but potentially complete sooner. For this to work properly, we must require something about the input too, namely that goals in the list are independent, that is, they must not communicate using shared variables or by any other means. From the caller's perspective,  \texttt{parallel/1} is a Prolog-style behavior: a pure-looking meta-predicate whose implementation happens to orchestrate actors.


Crucially, \texttt{parallel/1} does not introduce a new observable interaction pattern. Provided the listed goals are independent and free of observable interference, the call \texttt{parallel([G$_1$,\,..,\,G$_n$])} is extensionally equivalent to sequential conjunction for deterministic goals.

That is, it succeeds with exactly the same variable bindings, fails under the same conditions, and propagates the same errors. The only intended difference is cost. In this sense, \texttt{parallel/1} is an optimisation of sequential conjunction rather than a semantic extension. Its internal use of actors and asynchronous messaging remains observationally invisible, and it therefore remains within the sequential partition of the Prolog Web.



\subsubsection{Making RPC calls in parallel}\label{sec:parallel-rpc}

Two RPC calls in sequence take the running time of the first call plus the running time of the second:

\begin{verbatim}
?- rpc('http://n1.org', foo(X)),
   rpc('http://n2.org', bar(Y)).
\end{verbatim}

\noindent If the caller node has ACTOR capabilities, we may combine RPC calls with \texttt{parallel/1}:

\begin{verbatim}
?- parallel([
       rpc('http://n1.org', foo(X)),
       rpc('http://n2.org', bar(Y))
   ]).
\end{verbatim}

\noindent The time spent waiting is then approximately the maximum of the two response times. Note, however, that this is still a synchronous pattern: \texttt{parallel/1} blocks until both calls have returned. It therefore differs from actor-style interaction, where sends are asynchronous by default. 

In its basic form, \texttt{parallel/1} commits to the first solution of each goal; it does not explore combinations by backtracking, so if \texttt{foo/1} or \texttt{bar/1} are nondeterministic, the above returns only one solution from each call. To aggregate lists of solutions, use \texttt{findall/3}:

\begin{verbatim}
?- parallel([
       rpc('http://n1.org', findall(X, foo(X), Xs)),
       rpc('http://n2.org', findall(Y, bar(Y), Ys))
   ]),
   aggregate_solutions(Xs, Ys, Solutions).
\end{verbatim}

\subsubsection{Remote modules}\label{sec:remote-modules}

Below we sketch an implementation of \emph{remote modules}:\footnote{What we here call \emph{remote modules} is similar in spirit to Ciao Prolog's \emph{active modules}.}

\begin{verbatim}
use_remote_module(URI, ImportList, Options) :-
    maplist(import(URI, Options), ImportList).

import(URI, Options, Functor1/Arity as Functor2) :- !,
    functor(Head1, Functor1, Arity),
    Head1 =.. [Functor1|Args],
    Head2 =.. [Functor2|Args],
    assertz((Head2 :- rpc(URI, Head1, Options))).
import(URI, Options, Functor/Arity) :-
    functor(Head, Functor, Arity),
    assertz((Head :- rpc(URI, Head, Options))).
\end{verbatim}

\noindent Thanks to \texttt{load\_*} predicates, an import directive can define local wrappers even for predicates not directly exported by the remote node, provided they can be expressed in terms of predicates that are exported. 

\subsubsection{Ask around and seek agreement}\label{sec:other}

\begin{verbatim}
ask_around(URIs, Query) :-
    member(URI, URIs),
    catch(rpc(URI, Query), _, fail).

seek_agreement([], _Query).
seek_agreement([URI|URIs], Query) :-
    catch(rpc(URI, Query), _, fail),
    seek_agreement(URIs, Query).
\end{verbatim}



\subsubsection{Global goals}\label{sec:global-goals}

Our second example takes a step towards even more network transparency and is inspired by \cite{journals/jss/Loke06}. The idea is to support \emph{global goals} of the form \texttt{?- *Goal}, where a local peer database maps goal patterns to peers and options:

\begin{verbatim}
peer(foo(_), 'http://n1.org', [limit(10)]).
peer(bar(_), 'http://n2.org', [timeout(1),limit(100)]).
peer(bar(_), 'http://n3.org', [timeout(3),limit(100)]).
peer(baz(_), 'http://n3.org', [timeout(1),limit(50)]).
\end{verbatim}

\noindent Given such a database,\footnote{We make no assumptions on how the peer database is updated.} \texttt{*/1} can be implemented like so:

\begin{verbatim}
:- op(100, fx, *).

* Goal :-
    peer(Goal, Peer, Options),
    catch(rpc(Peer, Goal, Options), timelimit_exceeded, false).
\end{verbatim}

\noindent A registry-based variation can be obtained by retrieving \texttt{peer/3} information remotely:

\begin{verbatim}
* Goal :-
    rpc('http://registry.pl', peer(Goal, Peer, Options)),
    catch(rpc(Peer, Goal, Options), timelimit_exceeded, false).
\end{verbatim}


\subsection{Promise and yield over HTTP}
\label{sec:promise-yield}

The \texttt{rpc/2-3} predicate is synchronous: the caller waits until the remote node returns an answer. For many applications this is fine, but sometimes a client wants to initiate a remote computation, continue with other work, and collect the result later. The \texttt{promise/3-4} and \texttt{yield/2-3} predicates support this split-phase pattern.

A call to \texttt{promise/3-4} sends a query to a remote node but returns immediately, handing back a reference that identifies the pending computation:

\begin{verbatim}
?- promise('http://n1.org', mortal(Who), Ref).
Ref = 19303011.
?-
\end{verbatim}

\noindent The remote node begins evaluating \texttt{mortal(Who)}. Meanwhile, the caller is free to do other work. When it's ready for the answer, it calls \texttt{yield/2-3} with the reference:

\begin{verbatim}
?- yield($Ref, Answer).
Answer = success([mortal(plato), mortal(aristotle)], false).
?-
\end{verbatim}

\noindent If the answer has already arrived, \texttt{yield/2-3} returns immediately. If not, it waits. The answer is delivered to the caller's mailbox, so \texttt{yield/2-3} must be called by the same process that issued \texttt{promise/3-4}.

The familiar options carry over: \texttt{template/1} controls which bindings are returned (and can reduce payload size via template packing, cf.\ Section~\ref{sec:template-packing}), while \texttt{offset/1} and \texttt{limit/1} control slicing. By default, \texttt{promise/3-4} requests all solutions; with \texttt{limit(K)}, the answer contains at most $K$ solutions, allowing large result sets to be fetched in batches. It is often useful to think of these batches as \emph{pages} of an answer stream: \texttt{limit(K)} sets a page size, and \texttt{offset} selects which page to retrieve next. In Section~\ref{sec:backtracking-as-pagination} we will return to this idea and argue that pagination is best understood as a web-facing form of Prolog-style backtracking.

\subsubsection{Bounded waiting and timeouts}
\label{sec:bounded-waiting}

A client that cannot afford to block indefinitely can specify a timeout:

\begin{verbatim}
?- promise('http://n1.org', expensive_query(X), Ref),
   yield(Ref, Answer, [timeout(2.0)]).
?-
\end{verbatim}

\noindent If no answer arrives within two seconds, \texttt{yield/3} succeeds anyway -- but \texttt{Answer} remains unbound. The caller can test for this and decide how to proceed: retry, report an error, or fall back to a default.

A polling pattern emerges naturally. The following code initiates an asynchronous call and then periodically checks for the answer while doing other work:

\begin{verbatim}
?- promise('http://n1.org', mortal(Who), Ref, [limit(100)]),
   repeat,
     yield(Ref, Answer, [
       timeout(0.2),
       on_timeout(writeln('Still waiting...'))
     ]),
     (  var(Answer)
     -> do_other_work, fail
     ;  !
     ).
Still waiting...
Still waiting...
Answer = success([mortal(plato), mortal(aristotle)], false).
?-
\end{verbatim}

\noindent Each iteration waits briefly for an answer. On timeout, \texttt{on\_timeout/1} is called and \texttt{yield/3} succeeds without binding \texttt{Answer}, allowing the loop to continue. Once the answer arrives, the cut terminates the loop. 

This pattern -- non-blocking initiation followed by pull-based answer consumption -- preserves the sequential, declarative character of the Prolog Web while giving clients explicit control over timing.

\medskip
\noindent This is one instance of a more general rule: on the Prolog Web, timing belongs in the protocol, at the points where a process would otherwise block.

\section{The programmable Prolog Web}\label{sec:programmable-prolog-web}

The Prolog Web is programmable in a stronger sense than the familiar 'Web of APIs'. In its original form as a hypertext Web of
documents, the Web was not programmable at all. This changed early
with server-side scripts that generated documents dynamically, and
with client-side JavaScript. Later, remote procedure calling and
open web APIs made it natural to treat the Web as a platform for
composition, where independently developed services could be
combined into ``mashups'', often using a client language as glue.

Web Prolog supports this conventional style of integration, but also
generalises it. A node may offer clients not merely a fixed set of
endpoints, but a Web Prolog runtime environment that can be
programmed directly by queries of varying complexity and, when
permitted, by the injection of source code into long-lived processes
on the node. In addition, a node owner may install node-resident
predicates that become part of what clients can call. In this sense,
the Prolog Web is programmable both from the outside, by clients
issuing queries and spawning processes, and from the inside, by node
owners shaping the execution environment that clients interact with.

This dual perspective aligns naturally with the actor model. A node
may restrict which capabilities are present in an actor at birth,
for example via \texttt{load\_*} options that provide an initial
private program, and via placement choices (\texttt{node}) that
determine in which environment the actor should live. Some behaviour
is therefore ``innate'' to the actor, while other behaviour is
acquired through interaction. Both are subject to the policy,
resource limits, and trust assumptions of the hosting node.

Once code and data can cross node boundaries in either direction, a
central systems question arises. Should a computation be moved to
where the data resides, or should data (or code-as-data) be imported
to where the computation runs? This question has no single correct
answer. The optimal direction depends on data volume, update
frequency, reuse patterns, latency requirements, and trust. The
Prolog Web therefore does not enforce a fixed architecture, but
instead exposes a small set of orthogonal mechanisms that let
programmers choose, per interaction, how code and data should flow.

The owner's side of this exchange was illustrated in
Section~\ref{sec:node-resident-actors}: installing predicates in
the shared database and running node-resident actors that define
what clients can call. The remainder of this section focuses on the
client's side: how the \texttt{load\_*} options let a client
determine not only \emph{where} to run, but \emph{what} to run.


\subsection{Injecting code with the load\_* options}
\label{sec:injecting-code}

Looking back at the ping-pong scenario in
Section~\ref{sec:pong-pong-two-nodes}, what if the definition of
\texttt{pong/1} is available at \texttt{n1.org}, but not at
\texttt{n2.org}? The \texttt{load\_predicates} option ensures that
the appropriate source code is shipped from the caller's database
and injected into the private database of the actor that will run at
\texttt{n2.org}:

\begin{verbatim}
ping_pong :-
    spawn(pong, Pong_Pid, [
        node('http://n2.org'),
        load_predicates([pong/1])
    ]),
    spawn(ping(3, Pong_Pid)).
\end{verbatim}

\noindent The \texttt{load\_predicates} option names predicates that
already exist in the caller's context. When the clauses themselves
should be supplied literally, the \texttt{load\_list} option serves
that purpose. The following example uses it to populate a remote
toplevel actor's private database with a clause that participates in
the query:

\begin{verbatim}
?- toplevel_spawn(Pid, [
       node('http://n1.org'),
       load_list([(mortal(Who):-human(Who))])
   ]).
Pid = 83110912@'http://n1.org'.
?-   
\end{verbatim}

\noindent The query \texttt{?-mortal(Who)} now executes on
\texttt{n1.org}, combining the injected rule with the node's own
\texttt{human/1} facts:

\begin{verbatim}
?- toplevel_call($Pid, mortal(Who), [
       template(Who)
   ]).
?- receive({A -> true}).
A = success(83110912@'http://n1.org',[plato,aristotle],false).
?-
\end{verbatim}


\subsection{Code shipping and data locality}
\label{sec:code-shipping}

On the Internet, the cost of moving code and data across node
boundaries is significant, yet unavoidable. The practical question
is therefore not whether movement happens, but in which direction it
should happen in order to reduce cost and improve responsiveness. In
some cases it is more efficient to bring a small computation to a
large or frequently updated dataset; in others, it is better to
import data once and reuse it locally.

The most direct way to bring code to the data in Web Prolog is to
inject a small, task-specific predicate into a process placed near
the data, execute it there, and return only the derived result. This
corresponds closely to the classical systems idea of \textit{function
shipping}: rather than transferring a large dataset across the
network, one ships a small computation to where the data resides and
retrieves only a compact summary.

Conversely, bringing data to the code can be expressed by running a
goal locally while importing a remote program into the local
process. This is where a unifying feature of the design becomes
visible: since the default value of the \texttt{node} option for
\texttt{spawn/1-3} and \texttt{toplevel\_spawn/1-2} is the
special-purpose URI \texttt{localhost}, the same \texttt{load\_*}
mechanisms used for remote execution also work for local code
composition. The \texttt{load\_*} options are not a ``remote''
feature; they are a general code-composition mechanism that happens
to work transparently across the network. In this way, Web Prolog
treats code and data symmetrically: facts and rules can be moved,
cached, and reused much like datasets, reflecting Prolog's natural
blurring of the distinction between code and data.

These examples suggest a useful symmetry. Web Prolog code can flow
from client to node or from node to client, and the direction can be
chosen declaratively by configuring the actor or toplevel used for
the interaction. In principle, this choice can even be made
programmatically at runtime, allowing applications to adapt their
strategy based on observed latency, data size, or load.


\subsection{A small design space}
\label{sec:design-space}

Taken together, the mechanisms described in this chapter form a small
but expressive design space structured around a handful of recurring
choices:

\begin{itemize}
\item \textbf{Interface}: \texttt{spawn}, \texttt{toplevel\_spawn},
  \texttt{rpc}, or \texttt{promise/yield}?
\item \textbf{Placement}: locally or remotely (\texttt{node})?
\item \textbf{Code origin}: from the caller's database, shipped
  literally, or already resident on the target node
  (\texttt{load\_*})?
\item \textbf{Interaction style}: pure query, interactive
  enumeration, or actor-style orchestration?
\item \textbf{Synchrony}: should the caller block, or initiate work
  and collect the result later?
\item \textbf{State}: stateless request--response, or
  session-oriented exchange?
\end{itemize}

\noindent The point is not that one combination is always best, but
that the Prolog Web becomes programmable in a richer sense by letting
developers decide where computation happens and how programs and data
flow, rather than hard-coding those decisions into the architecture.

Code shipping is also the point where portability stops being an
academic concern and becomes operational. If nodes can execute code
originating elsewhere, then even small semantic differences or
library mismatches become visible immediately. For this reason, the
Prolog Web treats mobile code as profile-scoped: a piece of code
should state, implicitly or explicitly, which profile it assumes, so
that execution is either well-defined or rejected early.

In the limiting case, when a program is self-contained -- all predicates it depends on are either included or guaranteed by the target profile -- relocation becomes immediate. A client can download a program from one node and begin executing it locally, or ship it to another node, without any adaptation step. We refer to this property as \emph{instant portability}. It is not a special mechanism but a direct consequence of the design: profile-scoped code, symmetric code flow, and the uniform treatment of local and remote execution combine to make relocation a deployment decision rather than a porting effort.


\section{Timing and failure boundaries on the Prolog Web}
\label{sec:timing-failure-boundaries}

\noindent Distributed execution differs from local execution in one
decisive respect: waiting becomes observable. A goal that would
either succeed or fail immediately in a local Prolog system may,
when executed across a network, block for an indeterminate period.
Timeouts therefore cannot be treated as mere transport parameters;
they define the points at which partial failure becomes visible to
the programmer.

In the Prolog Web architecture, two transports are used -- HTTP and
WebSocket -- but three Web APIs are exposed: a stateless HTTP API, a
semi-stateful HTTP API, and a stateful WebSocket API. Each API
introduces a different blocking surface. Timeouts attach to those
surfaces. Their meaning, and therefore their appropriate values,
depend on the interaction model.


\subsection{Three blocking surfaces}

\begin{description}

\item[\textbf{Stateless HTTP API -- request--response blocking.}]
Blocking occurs at the boundary of a single HTTP request. A client
submits a goal and waits for a response. The timeout bounds that
roundtrip. Because each call is isolated and resources are allocated
per request, timeouts are typically short and strictly bounded.
Pagination via \texttt{limit} and \texttt{offset} exists precisely
to control cost and latency.

\item[\textbf{Semi-stateful HTTP API -- split-phase waiting.}]
Work submission and result retrieval are separated. A client
initiates a task and later invokes \texttt{yield/3} with a
\texttt{timeout/1} option. The timeout now bounds waiting for a
previously initiated computation. Short polling intervals, retries,
and fallback actions (\texttt{on\_timeout/1}) become explicit
application logic. Session-level idle or running limits govern
resource reclamation independently of per-call waiting.

\item[\textbf{Stateful WebSocket API -- mailbox waiting.}]
Blocking occurs inside \texttt{receive/2}. The timeout bounds
participation in an interaction protocol rather than a single
computation. A mailbox wait may represent a reply, a coordination
step, or a conversational turn. Its value is therefore determined by
protocol design, not merely transport latency. Longer waits are
natural in sustained dialogues, provided recovery behaviour is
defined.

\end{description}


\subsection{A comparative example}

Consider the goal
\texttt{route(gothenburg,\,paris,\,Route,\,Cost)}, which computes a
travel route together with its cost.

\begin{description}

\item[\textbf{Stateless HTTP.}]
The client poses the goal and waits for answers. The timeout bounds
the entire roundtrip. If exceeded, the call fails from the client's
perspective. The interaction remains goal-centric: one request, one
response. Expensive computations must be constrained (e.g.\ by
limiting search) or retrieved in slices.

\item[\textbf{Semi-stateful HTTP.}]
The client submits the goal but separates submission from waiting.
The node may continue computing after the initial response. The
client later calls \texttt{yield/3} with a short timeout, explicitly
managing when and how long to wait. Waiting becomes programmable
rather than implicit.

\item[\textbf{WebSocket.}]
The client sends a message and waits in \texttt{receive/2}. The
timeout now bounds participation in a dialogue rather than completion
of a single call. On timeout, the client may retry, redirect,
cancel, or renegotiate. The logical goal is unchanged, but the
interaction context has shifted from evaluation to conversation.

\end{description}

\noindent The difference is therefore not in what is computed, but in
where waiting is exposed and how failure is interpreted.


\subsection{Late replies}

Timeouts do not cancel the world. A computation may complete after a
client has stopped waiting. The question of what happens to the
result -- whether it is discarded, logged, or delivered if the
session remains alive -- is a question about the \emph{semantics of
timing}: it determines what the client can observe and what
guarantees the API provides. The Prolog Web makes these semantics
explicit at the language level through options such as
\texttt{timeout} and \texttt{on\_timeout}, rather than hiding them
inside transport configuration.

The complementary question -- what happens to the \emph{resources}
that produced the late reply, and who is responsible for reclaiming
them -- belongs to lifetime discipline and is addressed in
Section~\ref{sec:containment}.


\subsection{Design principle}

A simple rule emerges:

\begin{quote}
Place timeouts at the point where a process would otherwise block,
and choose their values according to the interaction model -- not
merely according to transport defaults.
\end{quote}

\noindent Under this discipline, timing becomes part of the
architectural contract. The two transports -- HTTP and WebSocket --
remain implementation details, but the three Web APIs expose distinct
failure boundaries. Making those boundaries explicit is essential for
predictable behaviour on an open network.


\section{Security on the Prolog Web}\label{sec:security-prolog-web}

The Prolog Web is designed to be open: any client that knows a node's
URI can submit goals for evaluation. Openness of this kind raises an
obvious question: what makes it safe to let strangers run code on
your node?

The answer has two components. Three mechanisms --
\emph{inexpressibility}, \emph{invisibility}, and
\emph{containment} -- define the \emph{structural} security
posture: the shape of what client code can express, what it can
reach, and how long its effects persist. These mechanisms arise directly from the language design, and are developed in detail below. 

The relevance of these mechanisms varies with the node profile. For
profiles such as RELATION and ISOBASE, where interaction is
stateless or request--response over HTTP, the dominant concerns are
those of conventional web services: input validation,
resource control, authentication, and transport protection. The
full capability discipline underlying invisibility and containment
becomes essential only for profiles that support long-lived sessions
and bidirectional communication, such as ISOTOPE and ACTOR, where
process identity and message passing cross node boundaries. In the
remainder of this section, we focus on this richer setting, where
the full interplay of the three mechanisms becomes visible.

A complete security story also requires
\emph{authentication} (to ground the distinction between client and
owner), \emph{resource governance} (to prevent abuse within the
sandbox), and \emph{transport security} (to protect communication on
the wire). These are deployment concerns that Web Prolog inherits
from standard web infrastructure; they are discussed briefly in
Section~\ref{sec:deployment-concerns}.

The layered structure of these mechanisms is illustrated in
Figure~\ref{fig:web-prolog-security-trinity}. The three lower boxes represent the structural security mechanisms
provided by Web Prolog itself; the upper layers correspond to
operational and deployment concerns.

\begin{figure}[ht]
\centering
\setlength{\fboxsep}{6pt}

\begin{minipage}{0.9\linewidth}
\centering

\fbox{\begin{minipage}{0.75\linewidth}
\centering
\textbf{Deployment security} \\
TLS -- authentication -- reverse proxy -- network isolation
\end{minipage}}

\vspace{0.5em}

\fbox{\begin{minipage}{0.75\linewidth}
\centering
\textbf{Resource governance} \\
quotas -- time limits -- spawn limits -- admission control
\end{minipage}}

\vspace{0.5em}

\begin{tabular}{ccc}
\fbox{\begin{minipage}{0.27\linewidth}
\centering
\textbf{Inexpressibility} \\
\small sandbox: restricted language subset
\end{minipage}}
&
\fbox{\begin{minipage}{0.27\linewidth}
\centering
\textbf{Invisibility} \\
\small capability discipline and name discipline
\end{minipage}}
&
\fbox{\begin{minipage}{0.27\linewidth}
\centering
\textbf{Containment} \\
\small lifetime discipline for client actors
\end{minipage}}
\end{tabular}

\end{minipage}

\vspace{0.5em}\caption{Layered security model of a Web Prolog node. The language sandbox restricts what client programs can express;
capability discipline controls which actors and services they may
reach; and lifetime containment ensures that client-created
computation cannot persist beyond the session. Resource governance limits the amount of computation a client may perform, while deployment mechanisms secure communication and identity at the network level.}
\label{fig:web-prolog-security-trinity}
\end{figure}


\subsection{Inexpressibility: the sandbox}\label{sec:inexpressibility}

Web Prolog adopts the same basic strategy as browser-based
JavaScript: client-injected code runs in a constrained language
subset. The language deliberately excludes file I/O, direct OS
access, raw socket programming, and persistent storage. These
capabilities, if needed, are provided by the host platform and
mediated by the node owner through predicates and services that the
owner chooses to install -- they are not expressible by
client-injected code.

This is the price of allowing strangers to run code on your machine:
the code must be incapable of causing certain classes of trouble.
The restriction is enforced at the language level, not by runtime
checks on individual calls, which means that a conforming
implementation can reject dangerous code statically rather than
catching violations at execution time.

Even when an operation exists in the underlying system, it is
exposed to clients only through owner-controlled predicates and
services, guarded by authentication and authorisation. The capability policy can therefore be read as part of the node's
security boundary: it defines which actions are available to which
role, and thereby what the node counts as its public surface.


\subsection{Invisibility: name discipline}\label{sec:invisibility}

Joe Armstrong articulated the principle clearly: if we do not know
the name of a process we cannot interact with it in any way, and the
system is secure; once names become widely known, the system becomes
less secure. Web Prolog inherits this principle directly.

A pid in Web Prolog is an opaque capability handle. The system does
not provide any client-facing process-enumeration interface, and pids
are chosen from a sufficiently large, unpredictable space that
guessing a live pid is computationally infeasible. A client can
therefore interact only with actors whose pids have been explicitly
revealed to it, subject to the node's authorisation policy. A pid is
a necessary capability -- we cannot address what we cannot name --
but not always sufficient; nodes may still require authorisation for
operations performed via a known pid.

Discoverability creates a natural tension with name discipline. If
processes are secure because you cannot talk to what you cannot name,
then public services must still provide a way for clients to learn
which names are meant to be used. The safe resolution is to treat
discovery as an owner-curated \emph{publication mechanism}: discovery
should resolve \emph{published} names, not reveal the existence of
everything that happens to be running. One concrete realisation of
this idea -- a declarative service directory queryable via
\texttt{rpc/2-3} -- was given in
Section~\ref{sec:naming-discovery}.

\subsubsection{Reply channels as capabilities}

A subtle but important consequence of treating pids as capabilities
concerns the common request--reply idiom. In many actor systems, a
client includes its own pid in a request so that the provider can
send the reply. In a closed system this is usually harmless, but on
the open Prolog Web it grants more authority than intended: a pid is
not merely a return address, but a general communication capability.

Revealing a client pid therefore allows the provider not only to
reply, but to send arbitrary messages to that actor. This
creates a form of unintended capability delegation that may be
benign in cooperative settings, but becomes problematic across
trust boundaries.

A simple discipline avoids this form of capability leakage. Instead of exposing its own pid, a client should generate a fresh, request-local reply
endpoint and include that in the request. Such an endpoint may be
realised as a temporary proxy actor or equivalent restricted mechanism, and
may be discarded after the reply has been received. In this way,
the provider receives authority only to deliver the designated
reply, not to invoke the client generally. In capability terms, the reply channel should carry only reply authority, not general send authority.

This pattern refines the capability discipline: communication authority should be as narrow and short-lived as the protocol step that requires it. This refinement is particularly important in distributed interactions, where reply channels cross node boundaries and therefore cross trust domains.

\subsection{Containment: lifetime discipline}\label{sec:containment}

On an open network, a client may disconnect mid-session or vanish
without warning. If server-side actors created on that client's
behalf continue to run, they retain memory and send messages into
mailboxes that may never be read. Over time, such orphans accumulate
into unbounded state and potential security exposure.

The preventive measure is to make process lifetimes
\emph{containment-shaped}: every actor spawned by a client is
lifetime-dependent on its parent, so that parent termination
deterministically tears down the entire subtree. In Web Prolog,
\texttt{link(true)} is the default; creating an independent actor
requires an explicit \texttt{link(false)}, and that option is
reserved for node owners. A client
cannot detach work from its session.

Remote placement sharpens the point. A client may place a child on a
remote node via \texttt{node(URI)}, but only under lifetime
dependence. Allowing a remote detached child would turn
\texttt{spawn/3} into a deployment primitive and make orphan
prevention intractable under disconnects and partitions. For that
reason, client-side remote spawn rejects \texttt{link(false)}.

Containment addresses the \emph{resource} side of the late-reply
problem identified in Section~\ref{sec:timing-failure-boundaries}:
when a client stops waiting, the computation it initiated may still
be running. Lifetime discipline ensures that such processes are
reclaimed when the session ends, regardless of whether their results
were ever consumed. The \emph{observable semantics} of what the
client sees at the timeout boundary -- whether a late reply is
discarded, delivered, or never generated -- is a timing concern and
is treated in that section.

Should client-originated detached jobs be needed in the future, they
should be introduced as an explicit abstraction -- with quotas,
lifetime parameters, and auditable handles -- rather than as a
general relaxation of the containment default.

\subsection{Deployment concerns}\label{sec:deployment-concerns}

The three structural mechanisms above define the security
\emph{architecture} of the Prolog Web, but they do not by themselves
constitute a complete deployment story.

First, the node must establish \emph{who} the caller is. Without
authentication, the distinction between client and owner has no firm
operational basis. Second, the node must enforce \emph{resource
governance}. A client may remain entirely within the sandbox and yet
still exhaust CPU time, memory, process slots, or network bandwidth.
Third, communication between client and node must be protected in
transit. Without TLS or equivalent transport protection, a third
party may observe or tamper with messages, undermining
confidentiality and integrity.

Fourth, the runtime process itself must be constrained by the host environment. Even a correctly specified language sandbox offers limited assurance if the underlying process retains unrestricted filesystem or network access; host isolation -- through operating-system-level confinement, WebAssembly runtimes, or WASI-based embeddings -- removes entire classes of risk that the language boundary alone cannot prevent. Section~\ref{sec:implementation-security} develops the relationship between language sandboxing and host isolation in detail.

These concerns are not part of the Web Prolog language design as
such. They belong to deployment, and are naturally addressed using
standard web infrastructure and operational controls. A production
node therefore requires both structural security and secure
deployment.

\medskip
\noindent What remains after these mechanisms have done their work is the
node's public surface: the predicates it exports, the services it
publishes, and the policies under which they may be accessed. A
well-configured node should therefore present a small, explicit, and
auditable interface: everything else should be either inexpressible,
invisible, or contained.


%\section{Security on the Prolog Web -- OLD}\label{sec:security-prolog-web}
%
%The Prolog Web is designed to be open: any client that knows a node's
%URI can submit goals for evaluation. Openness of this kind raises an
%obvious question: what makes it safe to let strangers run code on
%your node?
%
%The answer has two parts. Three mechanisms --
%\emph{inexpressibility}, \emph{invisibility}, and
%\emph{containment} -- define the \emph{structural} security
%posture: the shape of what client code can express, what it can
%reach, and how long its effects persist. These are the distinctive
%contributions of the language design, and they are developed in
%detail below. But a complete security story also requires
%\emph{authentication} (to ground the distinction between client and
%owner), \emph{resource governance} (to prevent abuse within the
%sandbox), and \emph{transport security} (to protect communication on
%the wire). These are deployment concerns that Web Prolog inherits
%from standard web infrastructure; they are discussed briefly in
%Section~\ref{sec:deployment-concerns}.
%
%
%\subsection{Inexpressibility: the sandbox}\label{sec:inexpressibility}
%
%Web Prolog adopts the same basic strategy as browser-based
%JavaScript: client-injected code runs in a constrained language
%subset. The language deliberately excludes file I/O, direct OS
%access, raw socket programming, and persistent storage. These
%capabilities, if needed, are provided by the host platform and
%mediated by the node owner through predicates and services that the
%owner chooses to install -- they are not expressible by
%client-injected code.
%
%This is the price of allowing strangers to run code on your machine:
%the code must be incapable of causing certain classes of trouble.
%The restriction is enforced at the language level, not by runtime
%checks on individual calls, which means that a conforming
%implementation can reject dangerous code statically rather than
%catching violations at execution time.
%
%Even when an operation exists in the underlying system, it is
%exposed to clients only through owner-controlled predicates and
%services, guarded by authentication and authorisation. The capability policy can therefore be read as part of the node's
%security boundary: it defines which actions are available to which
%role, and thereby what the node counts as its public surface.
%
%
%\subsection{Invisibility: name discipline}\label{sec:invisibility}
%
%Joe Armstrong articulated the principle clearly: if we do not know
%the name of a process we cannot interact with it in any way, and the
%system is secure; once names become widely known, the system becomes
%less secure. Web Prolog inherits this principle directly.
%
%A pid in Web Prolog is an opaque capability handle. The system does
%not provide any client-facing process-enumeration interface, and pids
%are chosen from a sufficiently large, unpredictable space that
%guessing a live pid is computationally infeasible. A client can
%therefore interact only with actors whose pids have been explicitly
%revealed to it, subject to the node's authorisation policy. A pid is
%a necessary capability -- we cannot address what we cannot name --
%but not always sufficient; nodes may still require authorisation for
%operations performed via a known pid.
%
%Discoverability creates a natural tension with name discipline. If
%processes are secure because you cannot talk to what you cannot name,
%then public services must still provide a way for clients to learn
%which names are meant to be used. The safe resolution is to treat
%discovery as an owner-curated \emph{publication mechanism}: discovery
%should resolve \emph{published} names, not reveal the existence of
%everything that happens to be running. One concrete realisation of
%this idea -- a declarative service directory queryable via
%\texttt{rpc/2-3} -- was given in
%Section~\ref{sec:naming-discovery}.
%
%
%\subsection{Containment: lifetime discipline}\label{sec:containment}
%
%On an open network, a client may disconnect mid-session or vanish
%without warning. If server-side actors created on that client's
%behalf continue to run, they retain memory and send messages into
%mailboxes that may never be read. Over time, such orphans accumulate
%into unbounded state and potential security exposure.
%
%The preventive measure is to make process lifetimes
%\emph{containment-shaped}: every actor spawned by a client is
%lifetime-dependent on its parent, so that parent termination
%deterministically tears down the entire subtree. In Web Prolog,
%\texttt{link(true)} is the default; creating an independent actor
%requires an explicit \texttt{link(false)}, and that option is
%reserved for node owners. A client
%cannot detach work from its session.
%
%Remote placement sharpens the point. A client may place a child on a
%remote node via \texttt{node(URI)}, but only under lifetime
%dependence. Allowing a remote detached child would turn
%\texttt{spawn/3} into a deployment primitive and make orphan
%prevention intractable under disconnects and partitions. For that
%reason, client-side remote spawn rejects \texttt{link(false)}.
%
%Containment addresses the \emph{resource} side of the late-reply
%problem identified in Section~\ref{sec:timing-failure-boundaries}:
%when a client stops waiting, the computation it initiated may still
%be running. Lifetime discipline ensures that such processes are
%reclaimed when the session ends, regardless of whether their results
%were ever consumed. The \emph{observable semantics} of what the
%client sees at the timeout boundary -- whether a late reply is
%discarded, delivered, or never generated -- is a timing concern and
%is treated in that section.
%
%Should client-originated detached jobs be needed in the future, they
%should be introduced as an explicit abstraction -- with quotas,
%lifetime parameters, and auditable handles -- rather than as a
%general relaxation of the containment default.
%
%
%\subsection{Deployment concerns}\label{sec:deployment-concerns}
%
%The three structural mechanisms above define what client code can do
%on a node. They do not, by themselves, constitute a complete security
%story. Three further concerns must be addressed at the deployment
%level.
%
%\textbf{Authentication.} Before inexpressibility and containment can
%do their work, the node must establish \emph{who} the caller is --
%at minimum, whether it is a client or an owner. Without
%authentication, the role distinction on which the entire capability
%policy rests has no foundation. Web Prolog does not prescribe a
%specific authentication mechanism; standard web infrastructure (TLS
%client certificates, API keys, OAuth tokens) provides the necessary
%building blocks.
%
%\textbf{Resource governance.} A client can stay entirely within the
%sandbox -- no forbidden operations, no name guessing, everything
%properly linked -- and still exhaust CPU, memory, or process slots.
%Containment ensures cleanup \emph{after} a session ends, but it does
%not prevent resource starvation \emph{during} the session. Quotas,
%time limits, and admission control are therefore a separate and
%necessary layer of defence. The timeout mechanisms described in
%Section~\ref{sec:timing-failure-boundaries} serve double duty here:
%they are part of the timing semantics visible to clients, but they
%also function as resource limits enforced by the node.
%
%\textbf{Transport security.} Messages between client and node travel
%over a network. Without TLS or equivalent protection, a third party
%can observe or tamper with communication, undermining both
%confidentiality and integrity. This concern is orthogonal to the
%structural mechanisms: a perfectly sandboxed, name-disciplined,
%containment-enforcing node is still vulnerable if its transport is
%unprotected.
%
%These are not afterthoughts but prerequisites. The structural
%mechanisms define the security \emph{architecture}; authentication,
%resource governance, and transport security define the security
%\emph{deployment}. A production node requires both.
%
%
%\subsection{The public surface}\label{sec:public-surface}
%
%What remains after these mechanisms have done their work is the
%node's \emph{public surface}: the predicates it exports, the
%services it publishes, and the policies under which they may be
%accessed. This surface is the node owner's choice and
%responsibility. It is defined positively -- by what the owner
%installs and advertises -- rather than negatively, by what the
%system fails to prevent.
%
%A well-configured node therefore presents a small, auditable
%interface. Clients interact with exported predicates through
%\texttt{rpc/2-3}, with published actors through registered names,
%and with nothing else. The structural mechanisms ensure that
%everything outside this interface is either inexpressible, invisible,
%or contained; the deployment mechanisms ensure that access is
%authenticated, resource-bounded, and protected in transit.


\section{Summary}\label{sec:prolog-web-summary}

This chapter developed the Prolog Web as the third component of the
Prolog Trinity, extending Web Prolog from a language for single-node
programming into a language for distributed systems.

The \emph{concurrent Prolog Web} introduced network-transparent
actor programming: spawning processes on remote nodes, sending
messages across node boundaries, and managing lifetimes through links
and monitors. Local and remote message passing differ only in
latency, not in kind. Node-resident actors demonstrated how an owner
can install durable, shared services, and a naming and discovery
mechanism showed how such services can be advertised without
compromising name discipline.

The \emph{sequential Prolog Web} provided a disciplined
request--response abstraction. The key construct, \texttt{rpc/2-3},
realises nondeterministic remote procedure calls with familiar Prolog
semantics, implemented first over a remote toplevel actor and then,
more practically, over stateless HTTP. The split-phase
\texttt{promise/3-4} and \texttt{yield/2-3} predicates extend the
pattern to deferred interaction. Although the sequential layer is
built on top of the concurrent substrate, it presents a model in
which mailboxes and scheduling remain invisible.

The \emph{programmable Prolog Web} showed how code and data flow
between client and node. The \texttt{load\_*} options let clients
inject task-specific predicates into remote execution contexts, while
the same mechanisms work locally via the \texttt{localhost} URI,
revealing code composition as a general facility rather than a
remote-only feature.

A section on \emph{timing and failure boundaries} made explicit
where waiting becomes observable across the three Web APIs --
stateless HTTP, semi-stateful HTTP, and WebSocket -- and established
that timeouts belong at blocking surfaces, chosen according to the
interaction model rather than transport defaults.

Finally, a section on \emph{security} asked what makes it safe to
let strangers run code on a node, and answered in terms of three
structural mechanisms -- inexpressibility, invisibility, and
containment -- together with the deployment concerns (authentication,
resource governance, and transport security) that complete the
picture. What these mechanisms leave exposed is the node's
\emph{public surface}: the predicates it exports, the services it
publishes, and the policies under which they may be accessed.

The chapters that follow build on this foundation. Part~II develops
generic behaviours -- servers, supervisors, toplevels, and
statecharts -- that structure actor implementations, and later parts
return to the architectural and paradigmatic questions that these
technical layers raise.


